\documentclass[10pt,oneside,openany]{memoir}

\iffalse
%font style 1
\usepackage[p,osf]{scholax}
\usepackage{amsmath,amsthm,amssymb}
\usepackage[scaled=1.075,ncf,vvarbb]{newtxmath}
\fi

%font style 2
\usepackage[T1]{fontenc}
\usepackage{amsmath}
\usepackage[fixamsmath]{mathtools}  % Extension to amsmath
\usepackage{amsthm}
\usepackage{amssymb}
\renewcommand*{\mathbf}[1]{\varmathbb{#1}}
\usepackage{newpxtext}
\usepackage{eulerpx}
\usepackage{mathrsfs}

\usepackage{eucal}
\usepackage{datetime}
    \newdateformat{specialdate}{\THEYEAR\ \monthname\ \THEDAY}
\usepackage[margin=1.5in]{geometry}
\usepackage{fancyhdr}
    \pagestyle{fancy}
    \fancyhf{}
    \cfoot{\scriptsize\thepage}
    \fancyhead[R]{\scalebox{0.8}{\nouppercase{\rightmark}}}
    \fancyhead[L]{\scalebox{0.8}{\nouppercase{\leftmark}}}
\usepackage{keytheorems}
    \newkeytheoremstyle{def}{
    inherit-style=definition,
    spaceabove=1.5em,
    spacebelow=1.5em,
    postheadspace=0.5em,
    }
    \newkeytheoremstyle{thm}{
    spaceabove=1.5em,
    spacebelow=1.5em,
    postheadspace=0.5em,
    }
    \newkeytheorem{theorem}[
    name = Theorem,
    style=thm,
    numberwithin = section,
    ]
    \newkeytheorem{theorem*}[
    name=Theorem,
    style=thm,
    numbered=false
    ]
    \newkeytheorem{proposition}[
    name = Proposition,
    style=thm,
    numberwithin = section,
    sibling=theorem
    ]
    \newkeytheorem{lemma}[
    name = Lemma,
    style=thm,
    numberwithin = section,
    sibling=theorem
    ]
    \newkeytheorem{corollary}[
    name = Corollary,
    style=thm,
    numberwithin = section,
    sibling=theorem
    ]
    \newkeytheorem{definition}[
    name = Definition,
    style=def,
    numberwithin = section,
    style = definition
    ]
    \newkeytheorem{example}[
    name = Example,
    style=def,
    numberwithin = section,
    style = definition
    ]
    \newkeytheorem{remark}[
    name=Remark,
    style=remark, 
    numbered=false  
    ]
    \newkeytheorem{question}[
    name=Question,
    style=remark, 
    numbered=false  
    ]
    \newkeytheorem{exercise}[
    name = Exercise,
    style=thm
    ]
    \newkeytheorem{axiom}[name = Axiom]
\usepackage{enumitem}
\usepackage[utf8x]{inputenc}
\usepackage{tikz}
\usepackage{tikz-cd}
\usepackage{float}
\usetikzlibrary{patterns}
\usetikzlibrary{positioning,arrows.meta}
\usetikzlibrary{decorations.pathreplacing}
\usepackage{wasysym}
\usepackage{hyperref}
\usepackage{ulem}
\usepackage{pgfplots}
\usetikzlibrary{patterns}
\usepgfplotslibrary{fillbetween}
\pgfplotsset{compat=1.18}

\renewcommand{\headrulewidth}{0pt} % removes the top/header line
% (optional) if you also want no line above the footer:
\renewcommand{\footrulewidth}{0pt}
%%%%%%%%%%%%%%%%%%%%%%%%%%%%%%%%%%%%%%%%%%%%%%%%%%%%%%%%%%%%%
%%%%%%%%%%%%%%%%%%%%%%%%%%%%%%%%%%%%%%%%%%%%%%%%%%%%%%%%%%%%%
% ============================================================
% makros.tex — reorganized, full restored version
% ============================================================

% ============================================================
% Sha / Cyrillic support notes
% ============================================================

%to make the correct symbol for Sha
%\newcommand\cyr{%
%\renewcommand\rmdefault{wncyr}%
%\renewcommand\sfdefault{wncyss}%
%\renewcommand\encodingdefault{OT2}%
%\normalfont \selectfont} \DeclareTextFontCommand{\textcyr}{\cyr}

% ============================================================
% Math operators
% ============================================================

\DeclareMathOperator{\ab}{ab}
\DeclareMathOperator{\ad}{ad}
\DeclareMathOperator{\adj}{adj}
\DeclareMathOperator{\alg}{alg}
\DeclareMathOperator{\Alt}{Alt}
\DeclareMathOperator{\Ann}{Ann}
\DeclareMathOperator{\arith}{arith}
\DeclareMathOperator{\Aut}{Aut}
\DeclareMathOperator{\Be}{B}
\DeclareMathOperator{\Bd}{Bd}
\DeclareMathOperator{\card}{card}
\DeclareMathOperator{\Char}{char}
\DeclareMathOperator{\csp}{csp}
\DeclareMathOperator{\codim}{codim}
\DeclareMathOperator{\coker}{coker}
\DeclareMathOperator{\coh}{H}
\DeclareMathOperator{\compl}{compl}
\DeclareMathOperator{\conj}{conj}
\DeclareMathOperator{\cont}{cont}
\DeclareMathOperator{\Cov}{Cov}
\DeclareMathOperator{\crys}{crys}
\DeclareMathOperator{\Crys}{Crys}
\DeclareMathOperator{\cusp}{cusp}
\DeclareMathOperator{\diag}{diag}
\DeclareMathOperator{\diam}{diam}
\DeclareMathOperator{\Dom}{Dom}
\DeclareMathOperator{\disc}{disc}
\DeclareMathOperator{\dist}{dist}
\DeclareMathOperator{\Div}{div}
\DeclareMathOperator{\dR}{dR}
\DeclareMathOperator{\Eis}{Eis}
\DeclareMathOperator{\End}{End}
\DeclareMathOperator{\ev}{ev}
\DeclareMathOperator{\eval}{eval}
\DeclareMathOperator{\Eq}{Eq}
\DeclareMathOperator{\Ext}{Ext}
\DeclareMathOperator{\Fil}{Fil}
\DeclareMathOperator{\Fitt}{Fitt}
\DeclareMathOperator{\Frob}{Frob}
\DeclareMathOperator{\G}{G}
\DeclareMathOperator{\Gal}{Gal}
\DeclareMathOperator{\GL}{GL}
\DeclareMathOperator{\Gr}{Gr}
\DeclareMathOperator{\Graph}{Graph}
\DeclareMathOperator{\GSp}{GSp}
\DeclareMathOperator{\GUn}{GU}
\DeclareMathOperator{\Hom}{Hom}
\DeclareMathOperator{\id}{id}
\DeclareMathOperator{\Id}{Id}
\DeclareMathOperator{\Ik}{Ik}
\DeclareMathOperator{\IM}{Im}
\DeclareMathOperator{\Image}{im}
\DeclareMathOperator{\Ind}{Ind}
\DeclareMathOperator{\Inf}{inf}
\DeclareMathOperator{\ins}{ins}
\DeclareMathOperator{\inv}{inv}
\DeclareMathOperator{\Isom}{Isom}
\DeclareMathOperator{\J}{J}
\DeclareMathOperator{\Jac}{Jac}
\DeclareMathOperator{\lcm}{lcm}
\DeclareMathOperator{\length}{length}
\DeclareMathOperator*{\limit}{limit}
\DeclareMathOperator{\Log}{Log}
\DeclareMathOperator{\M}{M}
\DeclareMathOperator{\Mat}{Mat}
\DeclareMathOperator{\N}{N}
\DeclareMathOperator{\Nm}{Nm}
\DeclareMathOperator{\NIk}{N-Ik}
\DeclareMathOperator{\NSK}{N-SK}
\DeclareMathOperator{\new}{new}
\DeclareMathOperator{\obj}{obj}
\DeclareMathOperator{\old}{old}
\DeclareMathOperator{\ord}{ord}
\DeclareMathOperator{\Or}{O}
\DeclareMathOperator{\op}{op}
\DeclareMathOperator{\out}{out}
\DeclareMathOperator{\PGL}{PGL}
\DeclareMathOperator{\PGSp}{PGSp}
\DeclareMathOperator{\rank}{rank}
\DeclareMathOperator{\Ran}{Ran}
\DeclareMathOperator{\rel}{Rel}
\DeclareMathOperator{\Real}{Re}
\DeclareMathOperator{\RES}{res}
\DeclareMathOperator{\Res}{Res}
%\DeclareMathOperator{\Sha}{\textcyr{Sh}}
\DeclareMathOperator{\Sel}{Sel}
\DeclareMathOperator{\semi}{ss}
\DeclareMathOperator{\sgn}{sign}
\DeclareMathOperator{\SK}{SK}
\DeclareMathOperator{\SL}{SL}
\DeclareMathOperator{\SO}{SO}
\DeclareMathOperator{\Sp}{Sp}
\DeclareMathOperator{\Span}{span}
\DeclareMathOperator{\Spec}{Spec}
\DeclareMathOperator{\spin}{spin}
\DeclareMathOperator{\st}{st}
\DeclareMathOperator{\St}{St}
\DeclareMathOperator{\SUn}{SU}
\DeclareMathOperator{\supp}{supp}
\DeclareMathOperator{\Sup}{sup}
\DeclareMathOperator{\Sym}{Sym}
\DeclareMathOperator{\Tam}{Tam}
\DeclareMathOperator{\tors}{tors}
\DeclareMathOperator{\tr}{tr}
\DeclareMathOperator{\Tr}{Tr}
\DeclareMathOperator{\un}{un}
\DeclareMathOperator{\Un}{U}
\DeclareMathOperator{\val}{val}
\DeclareMathOperator{\vol}{vol}

% ============================================================
% General aliases and short macros
% ============================================================

\newcommand{\absgal}{\G_{\bbQ}}
\newcommand{\Loc}{L_{\operatorname{loc}}}
\renewcommand{\k}{\kappa}
\newcommand{\Ff}{F_{f}}
%\newcommand{\ts}{\,^{t}\!}
\newcommand{\lat}{\mathscr{L}}
\newcommand{\ov}[1]{\overline{#1}}
\newcommand{\up}{\upsilon}
\newcommand{\e}{\epsilon}
\newcommand{\tac}{\textasteriskcentered}

%mahesh macros
\newcommand{\tm}{\textrm}

\newcommand{\bone}{\mathbf{1}}
\newcommand{\sign}{\mathrm{sign}}
\newcommand{\eps}{\varepsilon}
\newcommand{\textui}[1]{\underline{\textit{#1}}}

%\newcommand{\ov}{\overline}
%\newcommand{\un}{\underline}
\newcommand{\fin}{\mathrm{fin}}
\newcommand{\nl}{\newline \mbox{}}
\newcommand{\h}[1]{\hspace{#1pt}}
\newcommand{\mtext}[1]{\hspace{6pt}\text{#1}\hspace{6pt}}
\newcommand{\lnorm}{\left\lVert}
\newcommand{\rnorm}{\right\rVert}
\newcommand{\ds}{\displaystyle}
\newcommand{\ts}{\textstyle}
\newcommand{\sfrac}[2]{{}^{#1}\mskip -5mu/\mskip -3mu_{#2}}

% ============================================================
% Category-style operators and contoured variants
% ============================================================

\DeclareMathOperator{\Sets}{S \mkern1.04mu e \mkern1.04mu t \mkern1.04mu s}
    \newcommand{\cSets}{\scalebox{1.02}{\contour{black}{$\Sets$}}}
    
\DeclareMathOperator{\Groups}{G \mkern1.04mu r \mkern1.04mu o \mkern1.04mu u \mkern1.04mu p \mkern1.04mu s}
    \newcommand{\cGroups}{\scalebox{1.02}{\contour{black}{$\Groups$}}}

\DeclareMathOperator{\TTop}{T \mkern1.04mu o \mkern1.04mu p}
    \newcommand{\cTop}{\scalebox{1.02}{\contour{black}{$\TTop$}}}

\DeclareMathOperator{\Htp}{H \mkern1.04mu t \mkern1.04mu p}
    \newcommand{\cHtp}{\scalebox{1.02}{\contour{black}{$\Htp$}}}

\DeclareMathOperator{\Mod}{M \mkern1.04mu o \mkern1.04mu d}
    \newcommand{\cMod}{\scalebox{1.02}{\contour{black}{$\Mod$}}}

\DeclareMathOperator{\Ab}{A \mkern1.04mu b}
    \newcommand{\cAb}{\scalebox{1.02}{\contour{black}{$\Ab$}}}

\DeclareMathOperator{\Rings}{R \mkern1.04mu i \mkern1.04mu n \mkern1.04mu g \mkern1.04mu s}
    \newcommand{\cRings}{\scalebox{1.02}{\contour{black}{$\Rings$}}}

\DeclareMathOperator{\ComRings}{C \mkern1.04mu o \mkern1.04mu m \mkern1.04mu R \mkern1.04mu i \mkern1.04mu n \mkern1.04mu g \mkern1.04mu s}
    \newcommand{\cComRings}{\scalebox{1.05}{\contour{black}{$\ComRings$}}}

\DeclareMathOperator{\hHom}{H \mkern1.04mu o \mkern1.04mu m}
    \newcommand{\cHom}{\scalebox{1.02}{\contour{black}{$\hHom$}}}

% ============================================================
% Mathcal
% ============================================================

\newcommand{\cA}{\mathcal{A}}
\newcommand{\cB}{\mathcal{B}}
\newcommand{\cC}{\mathcal{C}}
\newcommand{\cD}{\mathcal{D}}
\newcommand{\cE}{\mathcal{E}}
\newcommand{\cF}{\mathcal{F}}
\newcommand{\cG}{\mathcal{G}}
\newcommand{\cH}{\mathcal{H}}
\newcommand{\cI}{\mathcal{I}}
\newcommand{\cJ}{\mathcal{J}}
\newcommand{\cK}{\mathcal{K}}
\newcommand{\cL}{\mathcal{L}}
\newcommand{\cM}{\mathcal{M}}
\newcommand{\cN}{\mathcal{N}}
\newcommand{\cO}{\mathcal{O}}
\newcommand{\cP}{\mathcal{P}}
\newcommand{\cQ}{\mathcal{Q}}
\newcommand{\cR}{\mathcal{R}}
\newcommand{\cS}{\mathcal{S}}
\newcommand{\cT}{\mathcal{T}}
\newcommand{\cU}{\mathcal{U}}
\newcommand{\cV}{\mathcal{V}}
\newcommand{\cW}{\mathcal{W}}
\newcommand{\cX}{\mathcal{X}}
\newcommand{\cY}{\mathcal{Y}}
\newcommand{\cZ}{\mathcal{Z}}

% ============================================================
% Mathscrl% 
%============================================================

\newcommand{\scra}{\mathscr{a}}
\newcommand{\scrA}{\mathscr{A}}
\newcommand{\scrb}{\mathscr{b}}
\newcommand{\scrB}{\mathscr{B}}
\newcommand{\scrc}{\mathscr{c}}
\newcommand{\scrC}{\mathscr{C}}
\newcommand{\scrd}{\mathscr{d}}
\newcommand{\scrD}{\mathscr{D}}
\newcommand{\scre}{\mathscr{e}}
\newcommand{\scrE}{\mathscr{E}}
\newcommand{\scrf}{\mathscr{f}}
\newcommand{\scrF}{\mathscr{F}}
\newcommand{\scrg}{\mathscr{g}}
\newcommand{\scrG}{\mathscr{G}}
\newcommand{\scrh}{\mathscr{h}}
\newcommand{\scrH}{\mathscr{H}}
\newcommand{\scri}{\mathscr{i}}
\newcommand{\scrI}{\mathscr{I}}
\newcommand{\scrj}{\mathscr{j}}
\newcommand{\scrJ}{\mathscr{J}}
\newcommand{\scrk}{\mathscr{k}}
\newcommand{\scrK}{\mathscr{K}}
\newcommand{\scrl}{\mathscr{l}}
\newcommand{\scrL}{\mathscr{L}}
\newcommand{\scrm}{\mathscr{m}}
\newcommand{\scrM}{\mathscr{M}}
\newcommand{\scrn}{\mathscr{n}}
\newcommand{\scrN}{\mathscr{N}}
\newcommand{\scro}{\mathscr{o}}
\newcommand{\scrO}{\mathscr{O}}
\newcommand{\scrp}{\mathscr{p}}
\newcommand{\scrP}{\mathscr{P}}
\newcommand{\scrq}{\mathscr{q}}
\newcommand{\scrQ}{\mathscr{Q}}
\newcommand{\scrr}{\mathscr{r}}
\newcommand{\scrR}{\mathscr{R}}
\newcommand{\scrs}{\mathscr{s}}
\newcommand{\scrS}{\mathscr{S}}
\newcommand{\scrt}{\mathscr{t}}
\newcommand{\scrT}{\mathscr{T}}
\newcommand{\scru}{\mathscr{u}}
\newcommand{\scrU}{\mathscr{U}}
\newcommand{\scrv}{\mathscr{v}}
\newcommand{\scrV}{\mathscr{V}}
\newcommand{\scrw}{\mathscr{w}}
\newcommand{\scrW}{\mathscr{W}}
\newcommand{\scrx}{\mathscr{x}}
\newcommand{\scrX}{\mathscr{X}}
\newcommand{\scry}{\mathscr{y}}
\newcommand{\scrY}{\mathscr{Y}}
\newcommand{\scrz}{\mathscr{z}}
\newcommand{\scrZ}{\mathscr{Z}}

% ============================================================
% Mathfrak (missing \fi)
% ============================================================

\newcommand{\fa}{\mathfrak{a}}
\newcommand{\fA}{\mathfrak{A}}
\newcommand{\fb}{\mathfrak{b}}
\newcommand{\fB}{\mathfrak{B}}
\newcommand{\fc}{\mathfrak{c}}
\newcommand{\fC}{\mathfrak{C}}
\newcommand{\fd}{\mathfrak{d}}
\newcommand{\fD}{\mathfrak{D}}
\newcommand{\fe}{\mathfrak{e}}
\newcommand{\fE}{\mathfrak{E}}
\newcommand{\ff}{\mathfrak{f}}
\newcommand{\fF}{\mathfrak{F}}
\newcommand{\fg}{\mathfrak{g}}
\newcommand{\fG}{\mathfrak{G}}
\newcommand{\fh}{\mathfrak{h}}
\newcommand{\fH}{\mathfrak{H}}
\newcommand{\fI}{\mathfrak{I}}
\newcommand{\fj}{\mathfrak{j}}
\newcommand{\fJ}{\mathfrak{J}}
\newcommand{\fk}{\mathfrak{k}}
\newcommand{\fK}{\mathfrak{K}}
\newcommand{\fl}{\mathfrak{l}}
\newcommand{\fL}{\mathfrak{L}}
\newcommand{\fm}{\mathfrak{m}}
\newcommand{\fM}{\mathfrak{M}}
\newcommand{\fn}{\mathfrak{n}}
\newcommand{\fN}{\mathfrak{N}}
\newcommand{\fo}{\mathfrak{o}}
\newcommand{\fO}{\mathfrak{O}}
\newcommand{\fp}{\mathfrak{p}}
\newcommand{\fP}{\mathfrak{P}}
\newcommand{\fq}{\mathfrak{q}}
\newcommand{\fQ}{\mathfrak{Q}}
\newcommand{\fr}{\mathfrak{r}}
\newcommand{\fR}{\mathfrak{R}}
\newcommand{\fs}{\mathfrak{s}}
\newcommand{\fS}{\mathfrak{S}}
\newcommand{\ft}{\mathfrak{t}}
\newcommand{\fT}{\mathfrak{T}}
\newcommand{\fu}{\mathfrak{u}}
\newcommand{\fU}{\mathfrak{U}}
\newcommand{\fv}{\mathfrak{v}}
\newcommand{\fV}{\mathfrak{V}}
\newcommand{\fw}{\mathfrak{w}}
\newcommand{\fW}{\mathfrak{W}}
\newcommand{\fx}{\mathfrak{x}}
\newcommand{\fX}{\mathfrak{X}}
\newcommand{\fy}{\mathfrak{y}}
\newcommand{\fY}{\mathfrak{Y}}
\newcommand{\fz}{\mathfrak{z}}
\newcommand{\fZ}{\mathfrak{Z}}

% ============================================================
% Mathbf
% ============================================================

\newcommand{\bfA}{\mathbf{A}}
\newcommand{\bfB}{\mathbf{B}}
\newcommand{\bfC}{\mathbf{C}}
\newcommand{\bfD}{\mathbf{D}}
\newcommand{\bfE}{\mathbf{E}}
\newcommand{\bfF}{\mathbf{F}}
\newcommand{\bfG}{\mathbf{G}}
\newcommand{\bfH}{\mathbf{H}}
\newcommand{\bfI}{\mathbf{I}}
\newcommand{\bfJ}{\mathbf{J}}
\newcommand{\bfK}{\mathbf{K}}
\newcommand{\bfL}{\mathbf{L}}
\newcommand{\bfM}{\mathbf{M}}
\newcommand{\bfN}{\mathbf{N}}
\newcommand{\bfO}{\mathbf{O}}
\newcommand{\bfP}{\mathbf{P}}
\newcommand{\bfQ}{\mathbf{Q}}
\newcommand{\bfR}{\mathbf{R}}
\newcommand{\bfS}{\mathbf{S}}
\newcommand{\bfT}{\mathbf{T}}
\newcommand{\bfU}{\mathbf{U}}
\newcommand{\bfV}{\mathbf{V}}
\newcommand{\bfW}{\mathbf{W}}
\newcommand{\bfX}{\mathbf{X}}
\newcommand{\bfY}{\mathbf{Y}}
\newcommand{\bfZ}{\mathbf{Z}}

\newcommand{\bfa}{\mathbf{a}}
\newcommand{\bfb}{\mathbf{b}}
\newcommand{\bfc}{\mathbf{c}}
\newcommand{\bfd}{\mathbf{d}}
\newcommand{\bfe}{\mathbf{e}}
\newcommand{\bff}{\mathbf{f}}
\newcommand{\bfg}{\mathbf{g}}
\newcommand{\bfh}{\mathbf{h}}
\newcommand{\bfi}{\mathbf{i}}
\newcommand{\bfj}{\mathbf{j}}
\newcommand{\bfk}{\mathbf{k}}
\newcommand{\bfl}{\mathbf{l}}
\newcommand{\bfm}{\mathbf{m}}
\newcommand{\bfn}{\mathbf{n}}
\newcommand{\bfo}{\mathbf{o}}
\newcommand{\bfp}{\mathbf{p}}
\newcommand{\bfq}{\mathbf{q}}
\newcommand{\bfr}{\mathbf{r}}
\newcommand{\bfs}{\mathbf{s}}
\newcommand{\bft}{\mathbf{t}}
\newcommand{\bfu}{\mathbf{u}}
\newcommand{\bfv}{\mathbf{v}}
\newcommand{\bfw}{\mathbf{w}}
\newcommand{\bfx}{\mathbf{x}}
\newcommand{\bfy}{\mathbf{y}}
\newcommand{\bfz}{\mathbf{z}}

% ============================================================
% Blackboard bold
% ============================================================

\newcommand{\bbA}{\mathbb{A}}
\newcommand{\bbB}{\mathbb{B}}
\newcommand{\bbC}{\mathbb{C}}
\newcommand{\bbD}{\mathbb{D}}
\newcommand{\bbE}{\mathbb{E}}
\newcommand{\bbF}{\mathbb{F}}
\newcommand{\bbG}{\mathbb{G}}
\newcommand{\bbH}{\mathbb{H}}
\newcommand{\bbI}{\mathbb{I}}
\newcommand{\bbJ}{\mathbb{J}}
\newcommand{\bbK}{\mathbb{K}}
\newcommand{\bbL}{\mathbb{L}}
\newcommand{\bbM}{\mathbb{M}}
\newcommand{\bbN}{\mathbb{N}}
\newcommand{\bbO}{\mathbb{O}}
\newcommand{\bbP}{\mathbb{P}}
\newcommand{\bbQ}{\mathbb{Q}}
\newcommand{\bbR}{\mathbb{R}}
\newcommand{\bbS}{\mathbb{S}}
\newcommand{\bbT}{\mathbb{T}}
\newcommand{\bbU}{\mathbb{U}}
\newcommand{\bbV}{\mathbb{V}}
\newcommand{\bbW}{\mathbb{W}}
\newcommand{\bbX}{\mathbb{X}}
\newcommand{\bbY}{\mathbb{Y}}
\newcommand{\bbZ}{\mathbb{Z}}

\newcommand{\Bmu}{\mbox{$\raisebox{-0.59ex}
  {$l$}\hspace{-0.18em}\mu\hspace{-0.88em}\raisebox{-0.98ex}{\scalebox{2}
  {$\color{white}.$}}\hspace{-0.416em}\raisebox{+0.88ex}
  {$\color{white}.$}\hspace{0.46em}$}{}}  %need graphicx and xcolor. this produces blackboard bold mu 

% ============================================================
% Greek-letter aliases
% ============================================================

\newcommand{\jota}{\jmath}

% ============================================================
% Matrices
% ============================================================

\newcommand{\bmat}{\left( \begin{matrix}}
\newcommand{\emat}{\end{matrix} \right)}

\newcommand{\bbmat}{\left[ \begin{matrix}}
\newcommand{\ebmat}{\end{matrix} \right]}

\newcommand{\pmat}{\left( \begin{smallmatrix}}
\newcommand{\epmat}{\end{smallmatrix} \right)}

\newcommand{\mat}[4]{\begin{pmatrix}{#1}&{#2}\\{#3}&{#4}\end{pmatrix}}

% ============================================================
% Residues and averaged integrals
% ============================================================

\newcommand{\res}[1]{\underset{#1}{\RES}\,}

\makeatletter
\newcommand*{\mint}[1]{%
  % #1: overlay symbol
  \mint@l{#1}{}%
}
\newcommand*{\mint@l}[2]{%
  % #1: overlay symbol
  % #2: limits
  \@ifnextchar\limits{%
    \mint@l{#1}%
  }{%
    \@ifnextchar\nolimits{%
      \mint@l{#1}%
    }{%
      \@ifnextchar\displaylimits{%
        \mint@l{#1}%
      }{%
        \mint@s{#2}{#1}%
      }%
    }%
  }%
}
\newcommand*{\mint@s}[2]{%
  % #1: limits
  % #2: overlay symbol
  \@ifnextchar_{%
    \mint@sub{#1}{#2}%
  }{%
    \@ifnextchar^{%
      \mint@sup{#1}{#2}%
    }{%
      \mint@{#1}{#2}{}{}%
    }%
  }%
}
\def\mint@sub#1#2_#3{%
  \@ifnextchar^{%
    \mint@sub@sup{#1}{#2}{#3}%
  }{%
    \mint@{#1}{#2}{#3}{}%
  }%
}
\def\mint@sup#1#2^#3{%
  \@ifnextchar_{%
    \mint@sup@sub{#1}{#2}{#3}%
  }{%
    \mint@{#1}{#2}{}{#3}%
  }%
}
\def\mint@sub@sup#1#2#3^#4{%
  \mint@{#1}{#2}{#3}{#4}%
}
\def\mint@sup@sub#1#2#3_#4{%
  \mint@{#1}{#2}{#4}{#3}%
}
\newcommand*{\mint@}[4]{%
  % #1: \limits, \nolimits, \displaylimits
  % #2: overlay symbol: -, =, ...
  % #3: subscript
  % #4: superscript
  \mathop{}%
  \mkern-\thinmuskip
  \mathchoice{%
    \mint@@{#1}{#2}{#3}{#4}%
        \displaystyle\textstyle\scriptstyle
  }{%
    \mint@@{#1}{#2}{#3}{#4}%
        \textstyle\scriptstyle\scriptstyle
  }{%
    \mint@@{#1}{#2}{#3}{#4}%
        \scriptstyle\scriptscriptstyle\scriptscriptstyle
  }{%
    \mint@@{#1}{#2}{#3}{#4}%
        \scriptscriptstyle\scriptscriptstyle\scriptscriptstyle
  }%
  \mkern-\thinmuskip
  \int#1%
  \ifx\\#3\\\else_{#3}\fi
  \ifx\\#4\\\else^{#4}\fi  
}
\newcommand*{\mint@@}[7]{%
  % #1: limits
  % #2: overlay symbol
  % #3: subscript
  % #4: superscript
  % #5: math style
  % #6: math style for overlay symbol
  % #7: math style for subscript/superscript
  \begingroup
    \sbox0{$#5\int\m@th$}%
    \sbox2{$#5\int_{}\m@th$}%
    \dimen2=\wd0 %
    % => \dimen2 = width of \int
    \let\mint@limits=#1\relax
    \ifx\mint@limits\relax
      \sbox4{$#5\int_{\kern1sp}^{\kern1sp}\m@th$}%
      \ifdim\wd4>\wd2 %
        \let\mint@limits=\nolimits
      \else
        \let\mint@limits=\limits
      \fi
    \fi
    \ifx\mint@limits\displaylimits
      \ifx#5\displaystyle
        \let\mint@limits=\limits
      \fi
    \fi
    \ifx\mint@limits\limits
      \sbox0{$#7#3\m@th$}%
      \sbox2{$#7#4\m@th$}%
      \ifdim\wd0>\dimen2 %
        \dimen2=\wd0 %
      \fi
      \ifdim\wd2>\dimen2 %
        \dimen2=\wd2 %
      \fi
    \fi
    \rlap{%
      $#5%
        \vcenter{%
          \hbox to\dimen2{%
            \hss
            $#6{#2}\m@th$%
            \hss
          }%
        }%
      $%
    }%
  \endgroup
}
\newcommand{\avg}{\mint{-}}
\makeatother

% ============================================================
% Comments and drafting helpers
% ============================================================

%Comments
\newcommand{\com}[1]{\vspace{5 mm}\par \noindent
\marginpar{\textsc{Comment}} \framebox{\begin{minipage}[c]{0.95
\textwidth} \tt #1 \end{minipage}}\vspace{5 mm}\par}

% ============================================================
% Arrows
% ============================================================

\newcommand{\hooktwoheadrightarrow}{%
  \hookrightarrow\mathrel{\mspace{-15mu}}\rightarrow
}

\makeatletter
\newcommand{\xhooktwoheadrightarrow}[2][]{%
  \lhook\joinrel
  \ext@arrow 0359\rightarrowfill@ {#1}{#2}%
  \mathrel{\mspace{-15mu}}\rightarrow
}
\makeatother

\makeatletter
\newcommand*{\da@rightarrow}{\mathchar"0\hexnumber@\symAMSa 4B }
\newcommand*{\da@leftarrow}{\mathchar"0\hexnumber@\symAMSa 4C }
\newcommand*{\xdashrightarrow}[2][]{%
  \mathrel{%
    \mathpalette{\da@xarrow{#1}{#2}{}\da@rightarrow{\,}{}}{}%
  }%
}
\newcommand{\xdashleftarrow}[2][]{%
  \mathrel{%
    \mathpalette{\da@xarrow{#1}{#2}\da@leftarrow{}{}{\,}}{}%
  }%
}
\newcommand*{\da@xarrow}[7]{%
  % #1: below
  % #2: above
  % #3: arrow left
  % #4: arrow right
  % #5: space left 
  % #6: space right
  % #7: math style 
  \sbox0{$\ifx#7\scriptstyle\scriptscriptstyle\else\scriptstyle\fi#5#1#6\m@th$}%
  \sbox2{$\ifx#7\scriptstyle\scriptscriptstyle\else\scriptstyle\fi#5#2#6\m@th$}%
  \sbox4{$#7\dabar@\m@th$}%
  \dimen@=\wd0 %
  \ifdim\wd2 >\dimen@
    \dimen@=\wd2 %   
  \fi
  \count@=2 %
  \def\da@bars{\dabar@\dabar@}%
  \@whiledim\count@\wd4<\dimen@\do{%
    \advance\count@\@ne
    \expandafter\def\expandafter\da@bars\expandafter{%
      \da@bars
      \dabar@ 
    }%
  }%  
  \mathrel{#3}%
  \mathrel{%   
    \mathop{\da@bars}\limits
    \ifx\\#1\\%
    \else
      _{\copy0}%
    \fi
    \ifx\\#2\\%
    \else
      ^{\copy2}%
    \fi
  }%   
  \mathrel{#4}%
}
\makeatother

% ============================================================
% Relation and delimiter spacing
% ============================================================

\renewcommand{\geq}{\geqslant}
\renewcommand{\leq}{\leqslant}
\newcommand{\midd}{\hspace{4pt}\middle|\hspace{4pt}}

% ============================================================
% Left Right Updated Deliminators
% ============================================================

\makeatletter
\NewDocumentCommand\Left{O{}}{%
    \IfValueTF{#1}{\notblank{#1}{\mathopen\bBigg@{#1}}{\left}}{\left}}
\NewDocumentCommand\Right{O{}}{%
    \IfValueTF{#1}{\notblank{#1}{\mathopen\bBigg@{#1}}{\right}}{\right}}
\makeatother

% ============================================================
% Restrictions
% ============================================================

\newcommand{\littletaller}{\mathchoice{\vphantom{\big|}}{}{}{}}
\newcommand\restr[2]{{% we make the whole thing an ordinary symbol
  \left.\kern-\nulldelimiterspace % automatically resize the bar with \right
  #1 % the function
  \littletaller % pretend it's a little taller at normal size
  \right|_{#2} % this is the delimiter
  }}

% ============================================================
% Chapter title formatting
% ============================================================

\newcommand{\chnum}{\titleformat
{\chapter} % command
[display] % shape
{\centering} % format
{\Huge \color{black} \shadowbox{\thechapter}} % label
{-0.5em} % sep (space between the number and title)
{\LARGE \color{black} \underline} % before-code
}

\newcommand{\chunnum}{\titleformat
{\chapter} % command
[display] % shape
{} % format
{} % label
{0em} % sep
{ \begin{flushright} \begin{tabular}{r}  \Huge \color{black}
} % before-code
[
\end{tabular} \end{flushright} \normalsize
] % after-code
}



\begin{document}
\begin{center}
    {\large Math 520 \\[0.1in]Week 6 Homework}\\[.175in]
    {Name:} {\underline{Gianluca Crescenzo\hspace*{2in}}}\\[0.15in]
    \end{center}
    \vspace{4pt}
%%%%%%%%%%%%%%%%%%%%%%%%%%%%%%%%%%%%%%%%%%%%%%%%%%
%%%%%%%%%%%%%%%%%%%%%%%%%%%%%%%%%%%%%%%%%%%%%%%%%%
%%%%%%%%%%%%%%%%%%%%%%%%%%%%%%%%%%%%%%%%%%%%%%%%%%
%%%%%%%%%%%%%%%%%%%%%%%%%%%%%%%%%%%%%%%%%%%%%%%%%%
\begin{exercise}
Let $H = L^2([0,1])$. Use the Gram-Schmidt orthogonalization process from linear algebra to produce an orthonormal set $\{u_1,u_2,u_3\}$ which spans the same space as $\Span\{1,x,x^2\}$.
\end{exercise}
\begin{proof}
    Let $v_1 := 1$, $v_2 := x$, and $v_3 := x^2$. The Gram-Schmidt process gives:
        \begin{equation*}
        \begin{split}
            r_1 &= v_1 = 1, \\
            \lnorm r_1 \rnorm &= \sqrt{(r_1,r_1)} = 1.
        \end{split}
        \end{equation*}

        \begin{equation*}
        \begin{split}
            r_2 &= v_2 - \frac{(v_2,r_1)}{(r_1,r_1)} r_1  = x - \frac{1}{2}. \\
            \lnorm r_2 \rnorm &= \sqrt{(r_2,r_2)} = \sqrt{\int_0^1 \left(x-\frac12\right)^2 dx} = \frac{1}{\sqrt{12}}
        \end{split}
        \end{equation*}

        \begin{equation*}
        \begin{split}
            r_3 &= v_3 - \frac{(v_3,r_1)}{(r_1,r_1)} r_1 - \frac{(v_3,r_2)}{(r_2,r_2)} r_2 = x^2 - x + \frac{1}{6}, \\
            \lnorm r_3 \rnorm & = \sqrt{(r_3,r_3)} = \sqrt{\int_0^1 \left(x^2-x+\frac16\right)^2 dx} = \frac{1}{\sqrt{180}}.
        \end{split}
        \end{equation*}

        Thus an orthonormal basis for $\Span\{1,x,x^2\}$ is:
            \begin{equation*}
            \begin{split}
                u_1 &:= \frac{r_1}{\lnorm r_1 \rnorm} = 1 \\
                u_2 &:= \frac{r_2}{\lnorm r_2 \rnorm} = \sqrt{12}\left(x-\frac{1}{2}\right) \\
                u_3 &:= \frac{r_3}{\lnorm r_3 \rnorm} = \sqrt{180} \left( x^2 - x + \frac{1}{6} \right). \qedhere
            \end{split}
            \end{equation*}
    \end{proof}

\newpage
\begin{exercise}
Show directly that the functions:
    \begin{equation*}
    \begin{split}
        u_n(x) = e^{i2\pi n x}\text{, }n \in \bfZ
    \end{split}
    \end{equation*}
form an orthonormal sequence in $L^1([0,1])$. Then write the function $f(x) = x$ in terms of this basis, that is find $a_n$ such that:
    \begin{equation*}
    \begin{split}
        x \sim \sum_{n = -\infty}^\infty a_n u_n(x).
    \end{split}
    \end{equation*}
What does Parseval's identity say in this case? (After some algebra you should be able to use your summation to find a formula for $\sum_{n = 1}^\infty \frac{1}{n^2}$.)
\end{exercise}
    \begin{proof}
        Observe that:
            \begin{equation*}
            \begin{split}
                \int_{0}^{1} e^{i 2 \pi n x}e^{-i 2 \pi n x}dx 
                & = \int_{0}^{1} 1 dx = 1. \\
                &\phantom{a} \\
                \int_{0}^{1} e^{i 2 \pi n x}e^{-i 2 \pi m x}dx
                & = \int_{0}^{1} e^{i 2 \pi (n-m)x}dx \\
                & = \frac{e^{i 2 \pi (n-m)x}}{i 2 \pi (n-m)}\Bigg|_0^1 \\
                & = \frac{e^{i 2 \pi (n-m)} - 1}{i 2 \pi (n-m)} \\
                & = \frac{\cos\bigl(2\pi(n-m)\bigr) - i\sin\bigl(2 \pi (n-m)\bigr) - 1}{i 2 \pi (n-m)} \\
                & = \frac{1 - 0 - 1}{i 2 \pi (n-m)} \\
                & = 0.
            \end{split}
            \end{equation*}
        Thus $\{u_n\}_{n = 1}^\infty$ is an orthonormal sequence in $L^1([0,1])$. We proceed by cases on computing the Fourier coefficients of $x$. For $n = 0$:
            \begin{equation*}
            \begin{split}
                a_0 = \int_0^1 x dx = \frac{1}{2}.
            \end{split}
            \end{equation*}
        For $n \neq 0$:
            \begin{equation*}
            \begin{split}
                a_n 
                & = (x,u_n) \\
                & = \int_0^1 x e^{-i 2 \pi n x}dx \\
                & = -\frac{e^{i 2 \pi n}(e^{i 2 \pi n} - i 2 \pi n - 1)}{4 n^2 \pi ^2} \h9\h9\text{\tiny From Mathematica} \\
                & = \frac{i}{2n \pi}. \h9\h9\text{\tiny Also from Mathematica}
            \end{split}
            \end{equation*}
        Using Parseval's Identity:
            \begin{equation*}
            \begin{split}
                \frac{1}{3}
                & = \int_0^1 x^2 dx \\
                & = \lnorm x \rnorm^2 \\
                & = \sum_{n = -\infty}^\infty |a_n|^2 \\
                & =|a_0|^2 + 2 \sum_{n = 1}^\infty |a_n|^2 \\
                & = \frac{1}{4} + 2 \sum_{n = 1}^\infty \left| \frac{i}{2 \pi n} \right|^2 \\
                & = \frac{1}{4} + \frac{1}{2\pi^2}\sum_{n = 1}^\infty \frac{1}{n^2}.
            \end{split}
            \end{equation*}
        Solving for $\sum_{n = 1}^\infty \frac{1}{n^2}$ gives:
            \begin{equation*}
            \begin{split}
                \sum_{n = 1}^\infty \frac{1}{n^2}
                & = 2\pi^2 \left( \frac{1}{3} - \frac{1}{4} \right) \\
                & = \frac{2\pi^2}{12} \\
                & = \frac{\pi^2}{6}. \qedhere
            \end{split}
            \end{equation*}
    \end{proof}

\begin{exercise}
Show directly that the functions:
    \begin{equation*}
    \begin{split}
        1, \cos x, \sin x, \cos 2x, \sin 2x, \cos 3x, ...
    \end{split}
    \end{equation*}
are mutually orthogonal in $L^2([-\pi,\pi])$.
\end{exercise}
\begin{proof}
    For $n,m \neq 0$:
        \begin{equation*}
        \begin{split}
            &\int_{-\pi}^{\pi}\cos (nx) dx = \frac{1}{n}\sin(nx)\h2\Bigg|_{-\pi}^{\pi} = \frac{1}{n}\sin(n \pi) - \frac{1}{n}\sin(-n\pi) = \frac{2}{n}\sin(n \pi) = 0,\\
            &\int_{-\pi}^{\pi}\sin (nx) dx = -\frac{1}{n}\cos(nx) \h2\Bigg|_{-\pi}^{\pi} = -\frac{1}{n}\cos(nx) + \frac{1}{n}\cos(-n\pi) = 0,\\
            &\int_{-\pi}^{\pi}\cos (nx) \cos (mx) dx = \int_{-\pi}^\pi \frac{\cos((n+m)x) + \cos((n-m)x)}{2}dx = 0,\\
            &\int_{-\pi}^{\pi}\sin (nx) \sin (mx) dx = \int_{-\pi}^\pi \frac{\cos((n-m)x) - \cos((n+m)x)}{2}dx = 0,\\
            &\int_{-\pi}^{\pi}\cos (nx) \sin (mx) dx= \int_{-\pi}^\pi \frac{\sin((n+m)x) + \sin((n - m)x)}{2}dx = 0. \qedhere
        \end{split}
        \end{equation*}
\end{proof}

\begin{exercise}
Let $\{u_n\}_{n = 1}^\infty$ be an orthonormal basis for a Hilbert space $H$. Let $\{v_n\}_{n = 1}^\infty$ be an orthonormal seq which satisfies $\sum_{n = 1}^\infty \lnorm u_n - v_n \rnorm^2 < 1$. Show that $\{v_n\}_{n = 1}^\infty$ is also a Hilbert basis for $H$.
\end{exercise}
    \begin{proof}
    Let $x \in \cH$. Suppose $x \perp \{v_j\}_{j = 1}^\infty$. By Parseval's Identity:
        \begin{equation*}
        \begin{split}
            \lnorm x \rnorm^2 
            & = \sum_{j = 1}^\infty |(x,u_n)|^2 \\
            & = \sum_{j = 1}^\infty |(x,u_n) - (x,v_n)|^2 \\
            & = \sum_{j = 1}^\infty |(x,u_n - v_n)|^2 \\
            & \leq \sum_{j = 1}^\infty \lnorm x \rnorm^2 \lnorm u_n - v_n \rnorm^2.
        \end{split}
        \end{equation*}
    Suppose towards contradiction $x \neq 0$. We can divide both sides of the above equation by $\lnorm x \rnorm^2$, giving:
        \begin{equation*}
        \begin{split}
            1 \leq \sum_{j = 1}^\infty \lnorm u_n - v_n \rnorm^2 < 1.
        \end{split}
        \end{equation*}
    This is a contradiction, hence $x = 0$. Since we've shown $x \perp \{v_j\}_{j = 1}^\infty$ implies $x = 0$, we've established $\{v_j\}_{j = 1}^\infty$ as an orthonormal basis for $\cH$.
    \end{proof}

\begin{exercise}\label{exercise:projection-norm-eq-1}
Let $H$ be a Hilbert space and $Y \neq \{0\}$ be a closed subspace, and let $P$ denote orthogonal projection onto $Y$. Show that:
    \begin{enumerate}[label = (\arabic*),itemsep=1pt,topsep=3pt]
        \item $P$ is linear;
        \item $P$ is bounded and $\lnorm P \rnorm = 1$;
        \item For all $x \in H$ we have $(Px,x) = \lnorm Px \rnorm^2$;
        \item $I - P$ is orthogonal projection onto $Y^\perp$.
    \end{enumerate}
\end{exercise}
\begin{proof}
    (a) For any $x_1,x_2 \in H$ and $\alpha \in \bfC$ we have:
        \begin{equation*}
        \begin{split}
            P(x_1 + \alpha x_2) 
            & = P(y_1 + y_1^\perp + \alpha y_2 + \alpha y_2^\perp) \\
            & = P(y_1 + \alpha y_2 + (y_1 + \alpha y_2)^\perp) \\
            & = y_1 + \alpha y_2 \\
            & = P(x_1) + \alpha P(x_2).
        \end{split}
        \end{equation*}
    Thus $T$ is linear.
    
    (b) Let $x \in H$. Observe that:
            \begin{equation*}
            \begin{split}
                \lnorm Px \rnorm^2 
                & = \lnorm y \rnorm^2 \\
                & \leq \lnorm y \rnorm^2 + \lnorm y^\perp \rnorm^2 \\
                & = \lnorm y + y^\perp \rnorm^2 \\
                & = \lnorm x \rnorm^2.
            \end{split}
            \end{equation*}
        Squaring both sides gives $\lnorm Px \rnorm \leq \lnorm x \rnorm$. Diving by $\lnorm x \rnorm$ on both sides of this inequality and taking the supremum over all $x \in H$, $x \neq 0$ gives $\lnorm P \rnorm_{\cL(H,Y)} \leq 1$. Now let $y \in B_Y(0,1)$. Then:
            \begin{equation*}
            \begin{split}
                \lnorm Py \rnorm = \lnorm y \rnorm = 1.
            \end{split}
            \end{equation*}
        Hence $1 \in \{\lnorm Px \rnorm \mid x \in B_H(0,1)\}$. It must be the case that $1 \leq \lnorm P \rnorm_{\cL(H,Y)}$. Therefore $\lnorm P \rnorm_{\cL(H,Y)} = 1$.

        (c) We will first show that $P$ is self-adjoint and idempotent. Let $x_1,x_2 \in H$. We have:
            \begin{equation*}
            \begin{split}
                (Px_1,x_2)
                & = (y,y_2 + y_2^\perp) \\
                & = (y_1,y_2) + (y_1,y_2^\perp) \\
                & = (y_1,y_2). 
            \end{split}
            \end{equation*}
            \begin{equation*}
            \begin{split}
                (x_1,Px_2)
                & = (y_1 + y_1^\perp , y_2) \\
                & = (y_1,y_2) + (y_1^\perp, y_2) \\
                & = (y_1,y_2).
            \end{split}
            \end{equation*}
        Therefore $P^\ast = P$; i.e., $P$ is self-adjoint. Now for any $x \in H$, notice that:
            \begin{equation*}
            \begin{split}
                P(P(x))
                & = P(y) \\
                & = y \\
                & = P(x).
            \end{split}
            \end{equation*}
        Thus $P^2 = P$. Using these two facts:
            \begin{equation*}
            \begin{split}
                \lnorm Px \rnorm^2 
                & = (Px,Px) \\
                & = (P^\ast P x, x) \\
                & = (P^2x,x) \\
                & = (Px,x). \qedhere
            \end{split}
            \end{equation*}

        (d) Let $y^\perp \in Y^\perp$. Then $(I - P)(y^\perp) = y^\perp$. Hence $(I-P)^2 = (I-P)$. Now for any $x \in H$: 
            \begin{equation*}
            \begin{split}
                (I - P)(x) 
                & = Ix - Px \\
                & = y+y^\perp - P(y + y^\perp) \\
                & = y+y^\perp - y \\
                & = y^\perp
            \end{split}
            \end{equation*}
        Hence $(I-P)$ is an orthogonal projection onto $Y^\perp$.
\end{proof}

\begin{exercise}
Let $H$ be a Hilbert Space with orthonormal basis $\{u_n\}_{n = 1}^\infty$. We try to define a map:
    \begin{equation*}
    \begin{split}
        T:H \rightarrow H \text{ by }Tx = \sum_{n = 1}^\infty \lambda_n (x,u_n)u_n
    \end{split}
    \end{equation*}
for a sequence $(\lambda_n) \subset \bfC$. First show that this map is a bounded linear operator if and only if $(\lambda_n)$ is a bounded sequence. Then, under what condition is $T$ self-adjoint?
\end{exercise}
\begin{proof}
    Suppose $T \in \cL(H)$. Then:
        \begin{equation*}
        \begin{split}
            |\lambda_n| 
            & = |\lambda_n| \lnorm u_n \rnorm \\
            & = \lnorm \lambda_n u_n \rnorm \\
            & = \lnorm T u_n \rnorm \\
            & \leq \lnorm T \rnorm \lnorm u_n \rnorm \\
            & = \lnorm T \rnorm.
        \end{split}
        \end{equation*}
    Taking the supremum over all $n \in \bfN$ gives $\sup_{n \in \bfN}|\lambda_n| \leq \lnorm T \rnorm_{\cL(H)} < \infty$. Therefore $(\lambda_n)$ is bounded. Conversely, suppose $\sup_{n \in \bfN}|\lambda_n| := M \leq \infty$. We must first show that $T$ is linear. Let $x,y \in H$ and $\alpha \in \bfC$. Define $a_n := (x,u_n)$ and $b_n := (y,u_n)$. We can write:
        \begin{equation*}
        \begin{split}
            x + \alpha y 
            & = \sum_{n = 1}^\infty a_n u_n + \alpha \sum_{n = 1}^\infty b_n u_n \\
            & = \sum_{n = 1}^\infty (a_n + \alpha b_n)u_n.
        \end{split}
        \end{equation*}
    Applying $T$ to $x+\alpha y$ gives:
        \begin{equation*}
        \begin{split}
            T(x+\alpha y) 
            & = \sum_{n = 1}^\infty \lambda_n (a_n + \alpha b_n)u_n \\
            & = \sum_{n = 1}^\infty \lambda_n a_n u_n + \alpha \sum_{n = 1}^\infty \lambda_n b_n u _n \\
            & = Tx + \alpha Ty.
        \end{split}
        \end{equation*}
    Thus $T$ is linear. Now let $x \in X$ and again define $a_n := (x,u_n)$. By Parseval's identity:
        \begin{equation*}
        \begin{split}
            \lnorm Tx \rnorm^2 
            & = \sum_{n = 1}^\infty \left|  \left( Tx, u_n \right) \right|^2 \\
            & = \sum_{n = 1}^\infty \left|  \left( \sum_{j = 1}^\infty \lambda_j a_n u_j, u_n \right) \right|^2 \\
            & = \sum_{n = 1}^\infty \left| \sum_{j = 1}^\infty \left( \lambda_j a_j u_j , u_n \right) \right|^2 \\
            & = \sum_{n = 1}^\infty \left| \lambda_n a_n \right|^2 \\
            & \leq \sum_{n = 1}^\infty |M|^2 |a_n|^2 \\
            & = |M|^2 \sum_{n = 1}^\infty |a_n|^2 \\
            & = |M|^2 \lnorm x \rnorm^2.
        \end{split}
        \end{equation*}
    Rearranging and taking the square-root of both sides yields:
        \begin{equation*}
        \begin{split}
            \frac{\lnorm Tx \rnorm}{\lnorm x \rnorm} \leq |M|.
        \end{split}
        \end{equation*}
    Taking the supremum over all $x \in X$, $x \neq 0$ gives $\lnorm T \rnorm_{\cL(H)} \leq |M| < \infty$. Thus $T \in \cL(H)$.

    Let $x,y \in H$. Observe that:
        \begin{equation*}
        \begin{split}
            \left( Tx, y \right)
            & = \left( \sum_{n = 1}^\infty \lambda_n (x,u_n)u_n, y \right) \\
            & = \sum_{n = 1}^\infty \bigl(\lambda_n (x,u_n)u_n,y\bigr) \\
            & = \sum_{n = 1}^\infty \lambda_n (x,u_n)(u_n,y).
        \end{split}
        \end{equation*}
        \begin{equation*}
        \begin{split}
            (x,Ty)
            & = \left( x,\sum_{n = 1}^\infty \lambda_n (y,u_n)u_n \right) \\
            & = \sum_{n = 1}^\infty \bigl(x, \lambda_n (y,u_n)u_n\bigr) \\
            & = \sum_{n = 1}^\infty \overline{\lambda_n (y,u_n)}(x,u_n) \\
            & = \sum_{n = 1}^\infty \overline{\lambda_n}(x,u_n)(y,u_n).
        \end{split}
        \end{equation*}
    From this, $T$ is self-adjoint if and only if $\lambda_n = \overline{\lambda_n}$; i.e., whenever $(\lambda_n) \subset \bfR$.
\end{proof}




\end{document}